\subsection{The conductor graph excludes Euler number one in the rational case}

\begin{statement}{P4.35}{Proved}
No irreducible support with rationally resolved normalization
has Euler number one. Consequently, in a fibration whose
supports are all irreducible or normal crossing, every
irreducible singular support with rationally resolved
normalization has Euler number exactly two.

More generally, an irreducible support cannot have reduced
image conductor equal to a rational curve with exactly one
point-identification cycle when its reduced normalization
conductor has no graph cycles. This excludes that conductor
case for every ruled-resolution double-horizontal pole tree.
\end{statement}
\begin{proof}
We first prove the graph assertion independently of
the transverse singularity type and of ramification.
Let $D\subset N$ be the reduced normalization conductor,
and suppose its point-identification graph is a tree.
Let $c$ be its number of normalized compact components.
Suppose its image $C$ is irreducible with rational
normalization and exactly one point-identification
cycle. There is then a unique point $p\in C$ with
two analytic branches; all other points are unibranch.
Denote the two preimages of $p$ on $C^\nu$ by
$p_+,p_-$. Every normalized component of $D$ maps
finitely and surjectively to $C^\nu$.

The finite set $\nu^{-1}(p)$ has $n$ points and
parametrizes the local irreducible hypersurface
branches of $S$ at $p$. Every pair of such branches
intersects along a curve: two nonunit local
hypersurface equations in a regular local ring of
dimension three generate an ideal of height at most
two. Their common curve lies in $C$. Consequently
the nonempty subsets of $\{+,-\}$ recording which
branches of $C$ are contained in each hypersurface
branch pairwise intersect. Such a family of subsets
of a two-element set has a common element. If a
singleton occurs its element belongs to every set;
if no singleton occurs every set is the whole set.
Interchanging the labels, every point of
$\nu^{-1}(p)$ therefore has at least one conductor
branch mapping to the plus branch of $C$.

Let $r_+$ and $r_-$ count the points on $D^\nu$
over $p_+$ and $p_-$ respectively. The preceding
paragraph gives $r_+\ge n$. Surjectivity of the
map from each of the $c$ normalized components
gives $r_-\ge c$. The contribution of the points
over $p$ to the point-identification graph of $D$
is exactly
\[
 \sum_{q\in\nu^{-1}(p)}(r_q-1)
       =r_++r_--n\ge c.
\]
For a connected reduced compact curve with $c$
normalized components, the cycle rank of this graph
is $1-c+\sum_q(r_q-1)$. All other contributions
are nonnegative. Its rank is therefore at least
one, contradicting the tree hypothesis.
This proves the graph assertion.

Now let the normalization have rational resolution
and suppose $e(S)=1$. The normal-crossing alternative
has Euler number two by P4.8. In the other alternative
P4.26 gives a connected simply connected resolved
pole support $A\subset T$, and
\[
 \pi_1(S)=\pi_1(C)=F_b,\qquad b\le1,
 \qquad e(C)=c_C+1-b.
\]
Its normalized components are rational.
Since $e(S_{\rm reg})\ge0$, Euler number one forces
$c_C=b=1$ and $e(S_{\rm reg})=0$.
It remains to verify the tree property for the
actual conductor $D\subset N$, allowing singular $N$.

P4.33 makes every singularity of $N$ rational.
Every connected exceptional fibre of $T\to N$
is simply connected. Here is the precise curve
calculation. Over a small Stein neighbourhood,
$R^1\rho_{N*}\mathcal O_T=0$ gives
$H^1(\mathcal O_T)=0$ on its inverse image.
This proper map has every fibre of dimension at
most one over that Stein neighbourhood. Its coherent
cohomological dimension is consequently at most
one, so the second higher direct image of every
coherent sheaf vanishes. Applied to the ideal of the reduced
exceptional curve $E$, Leray and the Stein property
give $H^2(\mathcal O_T(-E))=0$ there. Its ideal
sequence consequently gives $H^1(E,\mathcal O_E)=0$.
The curve normalization exact sequence now implies
that all normalized components are rational and
that their point-identification graph is a tree.
Thus $E$ is simply connected.

By P4R.6, $A$ is the full inverse image of $D$.
The proper map $A\to D$ is the quotient which
collapses these connected exceptional curves to
points and is an isomorphism elsewhere.
Compatible triangulation makes these curve
inclusions cofibrations. Collapsing a connected
simply connected subcomplex preserves fundamental
groups by van Kampen. Collapsing the finitely many
exceptional fibres therefore gives $\pi_1(D)=0$.
Each normalized component of $D$ is rational,
being birational to a strict pole component in
$T$. Hence its point-identification graph is a
tree. The already proved graph assertion contradicts
the description of $C$. This excludes $e(S)=1$.

For the final more general application, the
resolved pole support in the ruled double-section
case is $H$ with rational trees attached.
Every exceptional fibre is connected and is a
union of its components. If it does not contain
$H$, it is a connected rational subtree. If it
contains $H$, all the remaining components form
rational trees attached to that component once.
Contracting such a connected subcurve cannot
create a cycle in the graph: contraction of a
connected subgraph of a tree again has a tree
as quotient. The only positive-genus vertex,
if contracted, is replaced by a point, and this
also introduces no graph edge. Thus the image
conductor $D$ has no graph cycles, and the same
contradiction applies.

Finally, for a fibration in the two requested
support classes, P4.26 makes the total sum of
nonnegative fibre Euler numbers at most two.
The rational-resolution case is strictly positive;
excluding one leaves exactly two.
\end{proof}
